\section{Discussion}\label{sec:discussion}

\subsection{Strengths and weaknesses of the two approaches}\label{sec:strengths}

The comparison of Section~\ref{sec:llmresults} shows where a retrieval method trained for the task
and an LLM that reads the whole database each do well.

The proposed method is the more accurate of the two, and its lead grows with the size of the model:
it is smallest on the single-source cases and largest on the cross-source and synthesized cases,
whose ground truths combine sources or hold many equations (Table~\ref{tab:llmkind}). The method
reads the variable symbols of every equation and learns from cases with known ground truths how
the variables of a model fit together, so it can collect many equations through the variables they
share. It returns the same ranking for the same query, calls no external service, and trains in
minutes; the size of the database it can search is not bounded by an input limit. Its weaknesses
follow from the same reliance on variable overlap. It needs cases with known ground truths for
training, and it fills the top of the list with equations whose variables all belong to the query,
which pushes out ground-truth equations that share few variables with the query, such as
definitions and steady-state forms (Section~\ref{sec:erroranalysis}).

The LLM needs no training case. Without the variable lists and the context texts, Opus~5.5 placed
\RecAOpus{} of the ground-truth equations in its top 20 in Setting~A and beat the classical baseline
in every split. It found ground-truth equations that share few variables with the query, and on the cross-source cases its
MAP@20 hardly depended on the variable overlap; where it had the higher MAP@20, its first-ranked
equation was a ground-truth equation in \RankOneCrossOpus\% of its answers on the cross-source cases
and \RankOneSynthOpus\% on the synthesized cases. Its weaknesses are of four kinds. First, it is less accurate
overall, and the gap is widest on the models with several sources or many equations. Second, it costs money
at every query: one answer of Opus~5.5 cost \CostPerCallOpus{} US dollars on average,
\ReadShareOpus\% of it for reading the database, and this cost grows with the database and with the
number of queries. Third,
the input limit bounds what it can read: the database already takes \NDbTokens{} of the one million
tokens, which is why the variable definitions and the context texts stayed out, and a larger
database would not fit as it is. Fourth, its answers vary: between runs (the same list in
\SameListOpus\% of the cases), with the scale of the model (Haiku~5.5 fell below the classical
baseline), and in form (\ShareOpusAsInstructed\% of the answers of Opus~5.5 and
\ShareHaikuAsInstructed\% of those of Haiku~5.5 listed 20 identifiers as instructed).

These differences suggest different uses. The proposed method suits a large database, models that
draw on several sources or many equations, and repeated queries, provided that cases with known
ground truths exist. An LLM suits a field with no such cases: without any, Opus~5.5 beat the
classical baseline overall. In Recall@20, however, its advantage over the baseline came from the
synthesized cases alone; on the single-source cases, most of which ask for one equation, the
baseline was higher than Opus~5.5 in both metrics (Table~\ref{tab:llmkind}). The two approaches are also complementary. In the
cases in which Opus~5.5 had the higher Recall@20, the \NOnlyOpusEq{} ground-truth equations that
it alone placed (Section~\ref{sec:erroranalysis}) are of the kind that the proposed method tends to
push out, so an LLM that reranks or extends the candidates of the proposed method may recover such
equations at the cost of one short prompt, not of the whole database. We have not tested this
combination. Nor does the comparison show what an LLM could do with more information per equation:
we gave it the whole database at once, and the input limit left no room for the variable
definitions and the context texts that the proposed method reads.

\subsection{Matching variables by symbol}\label{sec:symbols}

The method matches variables by symbol alone and ignores their definitions. In \NSymbolCases{} of
the \NCases{} evaluation ground truths, a symbol that appears in two or more ground-truth equations
carries definitions that differ between them. Most differences are paraphrases, but a language
model (Claude Opus~5.5) that compared the definitions judged that in \NMeaningCases{} cases an input
or output symbol denotes different quantities in different equations.
We report no manual check of these judgments. On a sample of cases, a second language model, which
had not seen them, agreed with them in \NCheckFlagAgree{} of the \NCheckThreeApplicable{} cases that
hold a symbol in two or more ground-truth equations (\ref{app:check}).
Matching by symbol also ties the measured accuracy to the notation of the queries: every case takes
its input and output variables from the symbols of its ground-truth equations, so the experiments
cover only queries that share the notation of the ground truth. Because the six variable-overlap
features decide most of the accuracy (Section~\ref{sec:mechanism}), a query written in another
notation would lose the evidence on which the method relies. Judging whether two definitions name
the same quantity \citep{ProcessBERT} is the natural remedy and remains future work.

\subsection{Number of first-stage candidates}\label{sec:nonediscussion}

We chose $N_1$ as the smallest number of candidates, rounded up to the next thousand, that holds
every ground-truth equation of the training and validation cases. Under the weights
$(\WText, \WVar)$ this rule gave $N_1 = \NOne$ in every split, and the candidates then held every
ground-truth equation of the test cases as well. Over all \NCases{} cases, the first-stage ranking
places every ground-truth equation within its top \NOneAllGt; the top 50 hold \CovAFifty{} of the
ground-truth equations of a case on average, the top 200 hold \CovATwoHundred, and the top 500 hold
\CovAFiveHundred{} (\ref{app:n1}). A smaller $N_1$ loses the ground-truth equations that the first
stage ranks deepest, and a larger $N_1$ adds only incorrect candidates. Within the candidates, the
I/O Jaccard alone ranks a correct equation above an incorrect one drawn at random with probability
\AucCandIoJac{} (Table~\ref{tab:auc}), so incorrect candidates chosen at random are mostly easy,
which is consistent with hard negatives winning in every split. The rule does not ensure full
coverage on unseen cases, and the $N_1$ it yields may grow with the database.

\subsection{Training of sequential selection}\label{sec:tfdiscussion}

Learned greedy selection trains with teacher forcing: during training the reference set holds
ground-truth equations, whereas at inference it holds the model's own selections, which may include
mistakes. This gap may explain why learning adds little to greedy selection at inference
($\DiffMapAGreedyInfer$ and $\DiffMapBGreedyInfer$ in MAP). Training also adds the ground-truth
equations in one fixed order, that of their first-stage ranks. We compared no other order, such as
the order in which the equations determine one another's variables, and no remedy for the gap, such
as replacing some ground-truth equations of the reference set with the model's own selections
during training.

\subsection{Limitations}\label{sec:limitations}

Six limitations qualify the results. First, language models produced parts of the data: they
extracted the equations, wrote the \NEqRef{} equations of the generated equation sets, picked the ground
truths of the single-source cases, and wrote the descriptions of the synthesized cases and of part
of the cross-source cases, whose style may differ from requests written by engineers; a program
fixed the equations, inputs, and outputs of the synthesized cases. We report no manual verification
of this data. A language model, Claude Fable~5.2 \citep{ClaudeModels}, judged \NCheckCases{} cases,
drawn at random in proportion to the sizes of the three families, against four criteria
(\ref{app:check}). It judged that the ground-truth equations describe the model that the
description asks for in \NCheckOneYes{} cases, that the output variables are the quantities that
the description asks for in \NCheckTwoYes, that every symbol shared by two or more ground-truth
equations denotes one quantity in \NCheckThreeYes{} of the \NCheckThreeApplicable{} cases with such
a symbol, and that the equations determine the outputs once the inputs are given in
\NCheckFourYes. Only \NCheckAllMet{} cases met all four criteria, and none of the \NCheckProgram{}
cross-source cases that the program built did. The ground truths are therefore sets of equations
that the requirements of Section~\ref{sec:formulation} admit, not models that an engineer has
validated, and the reported accuracy measures how well a method recovers such sets.
Second, the evaluation is internal to the database: retrieval cannot find an equation that the
database lacks, and nobody screened the collected files by hand, so some have little to do with
physical models. Third, the queries share the notation of their ground truths
(Section~\ref{sec:symbols}) and, in \NReqThreeTrivial{} of the \NCases{} cases, list every variable of
the ground truth, which makes specificity decisive (Section~\ref{sec:mechanism}); only
\NWithIntermediate{} cases resemble the example of Fig.~\ref{fig:overview}, whose query omits the
intermediate variables of the model. Fourth, the split keeps the cases assigned to one source
together, but \ShareSharedEqA\% of the test cases of Setting~A share a ground-truth equation with a
training case (Section~\ref{sec:split}). On the test cases that share none, the proposed method
reached a Recall@20 of \RecANoShareProp{} and a MAP@20 of \MapTwANoShareProp{} and stayed above the
baseline and both LLMs in all 10 splits (\ref{app:overlap}); we did not retrain the reranker
without the training cases that share an equation with a test case. Fifth, among the hyperparameters, we compared no
alternative values of the reference-set size, the margin, the number of incorrect candidates, the
number of epochs, or the learning rate. Sixth, the comparison with LLMs covers two models of one
vendor, three runs per case, and one way of presenting the database.
